\begin{tabular}{lcr>{\raggedright\arraybackslash}p{5.4cm}}
\toprule
\textbf{Escenario} & \textbf{Rutas} & \textbf{Saltos} & \textbf{Motivo} \\
\midrule
Duitama -- Puente Nacional & 2 & 8 & caso central: dos rutas reales \\
Tunja -- Ubaté & 2 & 6 & ciclo: ambos extremos sobre el ciclo \\
Bucaramanga -- Bogotá & 2 & 6 & ciclo: lado de Santander contra Bogotá \\
Barbosa -- Zipaquirá & 2 & 6 & ciclo: origen en una hoja fuera del ciclo \\
Bogotá -- Villavicencio & 1 & 1 & ruta única, 1 salto \\
Girardot -- Cambao & 1 & 1 & ruta única, 1 salto, sin tocar Bogotá \\
Chiquinquirá -- Tunja & 1 & 2 & ruta única, 2 saltos \\
Bogotá -- Granada & 1 & 3 & ruta única, 3 saltos \\
Madrid -- Puerto Gaitán & 1 & 4 & ruta única, 4 saltos \\
Mariquita -- Granada & 1 & 6 & ruta única, 6 saltos (la más larga) \\
\bottomrule
\end{tabular}
